% ============================================================================
% Section IV: Solution Methods
% ============================================================================

\section{Solution Methods}
\label{sec:methods}

\gtopt{} implements three complementary solution methods, selected with
the \texttt{method} option.  All three solve the same formulation
(Section~\ref{sec:formulation}); they differ in how they decompose the
problem across time and scenarios.  The problem is organized in
\emph{scenes} (groups of scenarios) and \emph{phases} (groups of
stages); each (scene, phase) pair is an independent LP that is
assembled, solved and written in parallel.

\subsection{Monolithic LP}
\label{sec:monolithic}

The default method assembles the entire multi-scenario, multi-stage problem
into a single sparse LP/MIP and passes it to the solver backend.  The LP
is constructed in compressed sparse column (CSC) format via the
\texttt{FlatLinearProblem} class.

For a system with $N$ buses, $G$ generators, $D$ demands, $L$ lines,
$E$ batteries, $B$ blocks per stage, $T$ stages, and $S$ scenarios, the
approximate problem size is:

\begin{table}[ht]
  \centering
  \caption{LP Problem Size Estimates}
  \label{tab:lp_size}
  \begin{tabular}{@{}ll@{}}
    \toprule
    Quantity & Approximate count \\
    \midrule
    Operational variables & $(G{+}2D{+}2L{+}3E{+}N) \times S \times T \times B$ \\
    Investment variables  & $(G{+}D{+}L{+}E) \times T$ \\
    Bus balance rows      & $N \times S \times T \times B$ \\
    Kirchhoff rows        & $L \times S \times T \times B$ (if DC OPF) \\
    SoC balance rows      & $E \times S \times T \times B$ \\
    Capacity rows         & $(G{+}D{+}L{+}E) \times T$ \\
    \bottomrule
  \end{tabular}
\end{table}

The monolithic method is exact (no decomposition error) and is the
natural choice for daily and weekly operation, unit commitment, and
deterministic expansion problems that fit in memory.  It is also used,
with the SDDP cuts of a longer study as end-of-horizon valuation, to
solve the short-term scale of Section~\ref{sec:multiscale}.

\paragraph{Unit-commitment MIPs.}
For commitment problems, a feasible MIP start is built by a
relaxation-guided fix-and-propagate dive that satisfies the
exclusivity, minimum up/down and logic rows by construction; the start
is verified by the LP with all integers fixed, and, when it is
infeasible, an elastic LP identifies the integer variables to move
before a bounded sub-MIP is tried as a last resort.  Every MIP that
holds a feasible incumbent is then \emph{polished}: its integers are
fixed and the LP is re-solved, so that the reported objective, primal
solution and duals come from that LP, on every backend.  The time limit
is a single wall-clock budget per solve, shared by the start, the MIP
and the polish, and the prices are computed by the pricing modes of
Section~\ref{sec:sc}.  As an option, the pure commitment-logic rows
(transition, exclusion, startup tiers, charge/discharge exclusivity),
whose coefficients are all $\pm 1$, can be kept out of the row
equilibration, so that only the column scales apply to them.

\subsection{Stochastic Dual Dynamic Programming (SDDP)}
\label{sec:sddp}

For multi-stage stochastic problems, \gtopt{} implements SDDP following
Pereira and Pinto~\cite{pereira1991sddp}.  The phases are solved as
separate subproblems linked by state variables (reservoir volumes,
storage levels, allowance banks, irrigation-right balances) and by
Benders cuts that approximate the expected future cost
(Algorithm~\ref{alg:sddp}, Fig.~\ref{fig:sddp_flow}).

\begin{algorithm}[ht]
  \caption{SDDP Algorithm in \gtopt{}}
  \label{alg:sddp}
  \begin{algorithmic}[1]
    \STATE Initialize: empty cut sets $\mathcal{C}_t = \emptyset$
      for all $t$
    \FOR{iteration $k = 1, 2, \ldots, K$}
      \STATE \textbf{Forward pass}: Solve stages $t=1,\ldots,T$
        sequentially for sampled scenarios; record state variables
        (SoC, reservoir volumes)
      \STATE \textbf{Backward pass}: For $t=T,\ldots,1$, solve
        each stage for all aperture scenarios; compute dual
        variables; generate Benders cuts:
        \[
          \alpha_t \geq \hat{\alpha}_t + \hat{\pi}_t^\top
            (x_t - \hat{x}_t)
        \]
      \STATE Add cuts to $\mathcal{C}_t$
      \STATE Check convergence: bound gap $\leq \epsilon$, or a stationary gap
    \ENDFOR
  \end{algorithmic}
\end{algorithm}

\paragraph{Benders Optimality Cuts.}
The core mechanism of SDDP is the Benders optimality cut, which
approximates the expected cost-to-go function as a piecewise-linear
lower bound.  At each backward step for stage~$t$, the cut takes the
form:
\begin{equation}
  \label{eq:benders_cut}
  \alpha_t \;\geq\; \hat{z}_{t+1}
    + \sum_{i} \hat{\mu}_i \bigl(x_i - \hat{x}_i\bigr),
\end{equation}
where $\alpha_t$ is the future cost variable (the cost-to-go
approximation at stage~$t$), $\hat{z}_{t+1}$ is the optimal objective
value of the stage-$(t{+}1)$ subproblem evaluated at the trial point,
$\hat{\mu}_i$ are the dual multipliers of the state-linking constraints
(coupling successive stages through reservoir volumes, state-of-charge,
etc.), and $\hat{x}_i$ are the trial state variable values from the
forward pass.  As iterations proceed, the accumulated cuts progressively
tighten the piecewise-linear outer approximation of the true expected
cost-to-go, driving the lower bound upward toward the optimum.

\paragraph{Cut coefficients.}
The coefficients $\hat{\mu}_i$ in~\eqref{eq:benders_cut} are the
reduced costs of the dependent columns that receive the trial state in
stage~$t{+}1$, fixed at $\hat{x}_i$.  At an optimum these equal the duals
of explicit state-fixing rows, which therefore need not be added; the
LP is not modified between solves and its basis remains valid for warm
starts.  The coefficients are returned to physical
units with the scale factor of the column they are read from: under Ruiz
equilibration the state variable of stage~$t$ and its copy in
stage~$t{+}1$ are scaled independently, and mixing the two factors would
produce cuts that overestimate the future cost.

\paragraph{Uncertainty and apertures.}
Hydro inflows, renewable capacity factors and demand vary across
scenarios.  In the backward pass, each phase is re-solved for a set of
\emph{apertures}, i.e.\ alternative realizations of the uncertain data,
and the resulting cuts are averaged (single cut) or kept per aperture
(multi-cut).  Apertures can be sorted from wettest to driest and
truncated to the first $N$, which gives cheap early iterations and full
accuracy later; $N = 0$ gives the deterministic Benders backward pass
on the forward-pass solution.  Inflows can follow an AR(1) model, and
an experimental Markov-chain variant shares cuts by hydrological state.

The apertures of one (scene, phase) are grouped into parallel tasks in
one of three ways: \emph{per scene} (the default), where all apertures
are solved in series on one copy of the phase LP, each warm-started from
the optimal basis of the previous one; \emph{per aperture}, where each
aperture is a separate task with a fresh copy and a cold solve; and
\emph{chunked}, with $K$ apertures per task.  Since the apertures differ
only in column bounds, the per-scene grouping re-optimizes each of them
with a few dual-simplex pivots.  The grouping changes the scheduling and
the simplex path but not the cuts; it caps the backward-pass
parallelism at the number of scenes.  For backends without a warm start
(e.g.\ GPU first-order solvers), the cut of an aperture can instead be
synthesized from the duals of a representative aperture with an
intercept correction for the changed bounds, following Infanger and
Morton~\cite{infanger1996cutsharing}.

\paragraph{Feasibility.}
Hard limits such as minimum reservoir levels or irrigation deliveries
can make a forward-pass subproblem infeasible for a trial state.  An
elastic filter relaxes the state-linking bounds with penalized slacks,
identifies the violated limits, and adds a feasibility cut to the
previous phase, so that the policy learns to respect the limit instead
of paying a penalty.

In the backward pass, an aperture LP that is infeasible at the trial
state (or that fails to solve) is dropped, and the cut is averaged over
the remaining apertures with their probabilities renormalized.  This is
exact only under relatively complete recourse; otherwise, e.g.\ with a
hard end-of-horizon storage target that a dry aperture cannot reach,
the same aperture is dropped at every iteration and the lower bound
describes a rescaled, typically wetter, distribution.  Such iterations
are therefore never declared converged.  The drops are reported per
iteration, a persistent drop of the same (scene, phase, aperture) is
signalled once, and each run ends with aperture statistics (LPs solved,
dropped as infeasible or failed, persistent cases and the most affected
phases).  A soft terminal target restores relatively complete recourse.

\paragraph{Convergence.}
The upper bound is the mean cost of the forward-pass scenarios and the
lower bound is the root-phase objective with all cuts.  The algorithm
stops when the relative gap falls below a tolerance (default
$10^{-4}$), or, optionally, when the gap has been stationary over a
window of iterations, which accepts the irreducible gap of sampled
problems.  The stationary rule counts only the iterations of the
current solve (not bounds inherited from a previous cascade level) and
requires the gap to be below a ceiling (default 2\%).  Statistical
stopping rules based on the variance of the upper bound are also
available.  The lower bound of the future-cost variable $\alpha_t$ is
derived from the cuts actually installed, never imposed as
$\alpha_t \geq 0$, so models with negative future costs keep a valid
bound.  When the horizon is closed by external boundary cuts, the
expected boundary value is reported with the bounds and the gap is
corrected for it.

\paragraph{Engineering for large systems.}
Several features make SDDP practical on national-scale systems:
\begin{itemize}
  \item \textbf{Parallelism}: scenes, phases and apertures are solved in
    parallel by an adaptive work pool that throttles dispatch on free
    memory.
  \item \textbf{Low-memory mode}: each subproblem is kept as a
    compressed snapshot (zstd by default, lz4 optional) between solves
    and rebuilt on demand; the snapshots of the scenes of a phase are
    stored as compressed differences against one base and restored bit
    for bit.  Without compression, the solver backend of each
    subproblem is still released after its solve by default and rebuilt
    from an uncompressed snapshot with its saved basis, so that only the
    backends in flight are held in memory; keeping every LP live
    remains an option.
  \item \textbf{Warm starts}: simplex bases are stored in a
    solver-neutral form and reused across iterations, runs and
    backends; the forward solve of a phase can be seeded from the last
    basis of the backward aperture chain, which already prices the cuts
    just added.
  \item \textbf{Crossover policy}: with an interior-point forward
    solve, crossover is run only where a vertex is consumed, i.e.\ when
    the forward duals build the cuts (no apertures) or the forward
    basis is reused; otherwise the forward pass stops at the interior
    solution.  Cuts built from interior duals have their intercept
    lowered by the relative barrier gap, so that the primal--dual gap
    cannot raise the lower bound above the upper bound.
  \item \textbf{Cut management}: cuts are written as plain CSV after
    every iteration, so runs can be stopped, resumed (hot start) or
    extended; inactive cuts can be pruned and capped.
  \item \textbf{Boundary cuts}: an external future-cost function,
    e.g.\ the cuts of a longer study or of PLP, values the state at the
    end of the horizon.
\end{itemize}

\begin{figure}[t]
  \centering
  \resizebox{\columnwidth}{!}{%
  \begin{tikzpicture}[
    >=Stealth,
    stage/.style={draw, rounded corners=2pt, minimum width=0.95cm,
      minimum height=0.65cm, font=\scriptsize, thick},
    fwd/.style={stage, fill=blue!15},
    bwd/.style={stage, fill=orange!15},
    arr/.style={->, thick},
    stvar/.style={font=\tiny, blue!70},
    cutlbl/.style={font=\tiny, red!60},
  ]
    \node[font=\scriptsize\bfseries, blue!70, anchor=east] at (-0.7,1.8)
      {Forward};
    \node[font=\scriptsize\bfseries, orange!80!black, anchor=east]
      at (-0.7,0) {Backward};

    \foreach \n/\x/\l in {1/0/{t{=}1}, 2/1.9/{t{=}2}, T1/4.5/{T{-}1},
                          T/6.4/{t{=}T}} {
      \node[fwd] (f\n) at (\x,1.8) {$\l$};
      \node[bwd] (b\n) at (\x,0) {$\l$};
    }
    \node[font=\scriptsize] at (3.2,1.8) {$\cdots$};
    \node[font=\scriptsize] at (3.2,0) {$\cdots$};

    % forward: state passed to the next stage
    \draw[arr, blue!60] (f1) -- (f2) node[stvar, above, midway] {$x_1$};
    \draw[arr, blue!60] (f2) -- (2.85,1.8);
    \draw[arr, blue!60] (3.55,1.8) -- (fT1);
    \draw[arr, blue!60] (fT1) -- (fT)
      node[stvar, above, midway] {$x_{T-1}$};

    % backward: cuts passed to the previous stage
    \draw[arr, red!50] (bT) -- (bT1) node[cutlbl, below, midway] {cut};
    \draw[arr, red!50] (bT1) -- (3.55,0);
    \draw[arr, red!50] (2.85,0) -- (b2);
    \draw[arr, red!50] (b2) -- (b1) node[cutlbl, below, midway] {cut};

    % trial states handed from forward to backward pass
    \draw[->, dashed, gray, thin] (fT.south) -- (bT.north)
      node[font=\tiny, midway, left] {trial $\hat{x}$};

    % convergence check and loop
    \node[draw, rounded corners=3pt, fill=gray!10, font=\scriptsize]
      (conv) at (3.2,-1.3)
      {$|z^\text{UB} - z^\text{LB}|/z^\text{UB} \leq \varepsilon$?};
    \draw[->, thick, gray] (b1.south) |- (conv.west);
    \draw[->, thick, gray, dashed] (conv.east) -| (7.4,2.7) -| (f1.north);
    \node[font=\tiny, gray, anchor=west] at (7.45,0.9) {iterate};
  \end{tikzpicture}}%
  \caption{SDDP forward-backward iteration: the forward pass solves
    stages sequentially passing state variables (battery SoC, reservoir
    volume); the backward pass generates Benders cuts. Convergence is
    checked via the upper/lower bound gap.}
  \label{fig:sddp_flow}
\end{figure}

\subsection{Cascade Decomposition}
\label{sec:cascade}

The cascade method runs SDDP over a sequence of levels with
progressively more detailed models.  Each level defines its own LP
options (e.g.\ single bus, then the full network without losses, then
the network with losses) and SDDP options (iterations, apertures,
tolerances).  A level starts from the cuts of the previous level, kept
when their duals remain valid and otherwise re-generated, so the
expensive detailed model needs few iterations.  A typical configuration
uses 4 apertures in level~0, 8 in level~1 and the full set in level~2.
Each level may also override the solver options of its forward and
backward passes (tolerances, MIP gap, time limit, algorithm).

A level may also use its own, reduced network, produced by a network
reduction tool that merges buses and aggregates lines while keeping
every storage element.  Cuts are transferred between levels with a key
per coefficient (element class, variable and identifier) instead of a
column index, so they remain applicable when the LP is rebuilt with a
different formulation or network.  A reduced level whose network is a
relaxation of the next one (merged buses, summed inter-zonal capacities,
no Kirchhoff rows and no losses) has a cost-to-go no greater than that
of the detailed network, so its cuts are valid lower bounds and can be
kept indefinitely; cuts of a reduced level with equivalent reactances
or resistances are not guaranteed to be valid and are inherited for a
limited number of iterations.  Each level reports its bounds, gap,
iterations and aperture statistics.

\subsection{Investment Master}
\label{sec:investment_master}

For expansion problems that are too large for a monolithic LP, an
investment master in the style of OptGen chooses the investment
modules, and SDDP evaluates the expected operating cost of each plan.
Benders cuts on the investment variables return the marginal value of
each module to the master.

\subsection{One Model Across Time Scales}
\label{sec:multiscale}

Because the same element definitions are used at every scale, only the
time structure (stages and blocks of any duration) and the method change
between studies (Table~\ref{tab:scales}).

\begin{table}[t]
  \centering
  \caption{Time scales and methods}
  \label{tab:scales}
  \begin{tabular}{@{}llll@{}}
    \toprule
    Scale & Horizon & Blocks & Method \\
    \midrule
    Operation   & hours--days   & 15\,min--1\,h      & monolithic (UC) \\
    Medium term & months--years & hours, load curve & SDDP \\
    Planning    & up to 30+ yr  & representative    & SDDP, cascade, \\
                &               &                   & investment master \\
    \bottomrule
  \end{tabular}
\end{table}

The scales are coupled in two directions:
\begin{itemize}
  \item \textbf{Optimality cuts} carry water values and storage
    opportunity costs downward: the future-cost function of a planning
    or medium-term study becomes the boundary cut set of a shorter
    study, and the cascade inherits the cuts of its coarser levels.
  \item \textbf{Feasibility cuts} enforce reservoir levels and water
    rights at each scale, so the policy handed from one scale to the
    next is also feasible.
\end{itemize}

As a result, the water value used to dispatch a reservoir on a given day
is the same marginal value that the medium-term study computed for that
reservoir, and the investment plan is evaluated with the same operating
model that will later schedule it.

\paragraph{Typical periods in expansion studies.}
For expansion studies the year is usually reduced to one chronological
typical week (or a few representative days) per month, on which
storage cycles, ramps and commitment run in real time while energy,
costs and annuities are scaled to the month (Section~\ref{sec:point_stages}).
The converters provide several experimental week samplers: uniform
blocks of $H$ hours, finer blocks at the daily net-load peak and the
steepest ramps, and point samples at $N$ Chebyshev nodes with
quadrature weights.  The default window runs from Wednesday noon to
Wednesday noon, which places the weekend, where low load and high solar
drive storage cycling and curtailment, away from the cyclic closure, and
the week holding the monthly net-load peak (or the largest net energy) is
selected.  An optional refinement adds 30- or 15-minute resolution
around the sunrise and sunset ramps detected in the data: block samplers
split the base blocks there with energy-conserving, shape-preserving
cell means, and Chebyshev weeks add the nodes of the order-$2N$
Lobatto set nearest to each event, which contains the order-$N$ set,
with the weight of each dropped node split between its kept
neighbours so that all weights stay positive.  Commitment is not
refined, so the refinement adds continuous variables only.  A separate
converter, \texttt{tcheby2gtopt}, translates the validation systems of
the original Chebyshev scheduling code~\cite{matus2015chebyshev} into
both representations.

In a preliminary comparison on one expansion year of the Australian
National Electricity Market (twelve typical weeks with 138 batteries,
solved with HiGHS), the objective of the Chebyshev weeks approached the
reference with one-hour blocks from below as $N$ grew: 0.9 to 1.3\% below
it with 48 to 64 nodes and within 0.5\% with 84 nodes, against 0.24\%
for uniform two-hour blocks of the same LP size.  With 24 nodes the
week was under-resolved and lost about a third of the battery cycling.  The \texttt{cc} rule~\eqref{eq:point_cc} was closest to the
reference among the storage rules and gave the smallest LP, with the
same size as a block LP with the same number of blocks; the banded
rules add the envelope rows.  The converters therefore write
\texttt{cc} for every storage on point stages.

\subsection{Solver Backends}

All methods delegate LP/MIP solution to a plugin
(\texttt{libgtopt\_solver\_*.so}) loaded at run time.  The supported
backends are CLP/CBC~\cite{forrest2005cbc},
HiGHS~\cite{huangfu2018highs}, SCIP, NVIDIA cuOpt (GPU), CPLEX, Gurobi
and MindOpt, each declaring its capabilities (MIP, basis input and
output) to the core (Section~\ref{sec:architecture}).  Without an
explicit choice, the first available of CPLEX, HiGHS, MindOpt, CBC and
CLP is used.  Each solve reports its status and
an estimate of the condition number, and per-element variable scaling,
objective scaling and row/column equilibration (Ruiz by default for
multi-bus Kirchhoff networks, row-max otherwise) keep models with
kWh-scale batteries and TWh-scale reservoirs numerically stable.
